% !PRINTED 698
\weberpage{698}
If now $\gamma = 0$, it follows from this, on letting $t$ and $m$ grow to infinity, that $T : t$ too has the limiting value 0. But if $\gamma$ is different from 0, then by (7)
\[
\operatorname{Lim} \frac{\mu_m}{\mu_{m+1}} = \operatorname{Lim} \frac{m}{m+1} = 1,
\]
and therefore the two limiting values in (9) approach the value $\gamma$, and accordingly there follows:
\weq{10}{\mathop{\mathrm{Lim}}\limits_{t=\infty} \frac{T}{t} = \gamma.}

If, secondly, we conversely assume the limit equation (10), then here too $\mu_n$ will have to grow to infinity with $n$; for otherwise, since the total number of all $\mu$ is infinite, $T$ would already be infinite for a finite $t$, and $\gamma$ could not be finite.

Let us now take a value $\mu_n$ which occurs in the sequence of the mutually equal
\[
\mu_{m+1}, \; \mu_{m+2}, \; \dots, \; \mu_{m+l}
\]
so that
\[
m + 1 \leqq n \leqq m + l,
\]
then $T$, when $t$ passes through the value $\mu_n$, will suddenly increase by $l$ units, and the ratio $T : t$ increases by $l : \mu_n$. But since $T : t$ is to have a finite limiting value, we must have
\weq{11}{\operatorname{Lim} \frac{l}{\mu_n} = 0.}
Now if $t = \mu_n$, then $T = m + l$, and consequently
\weq{12}{\frac{T}{t} = \frac{m+l}{\mu_n}, \quad \operatorname{Lim} \frac{m+l}{\mu_n} = \gamma.}

Further
\[
\frac{m}{\mu_n} < \frac{n}{\mu_n} \leqq \frac{m+l}{\mu_n},
\]
and consequently, by (11) and (12),
\[
\operatorname{Lim} \frac{n}{\mu_n} = \operatorname{Lim} \frac{m}{\mu_n} = \gamma.
\]
Thus equation (7) has been deduced from equation (10) \textsuperscript{1}).

By using theorem 4. we can now also give theorem 3. the form:

\begingroup
\footnotesize
\textsuperscript{1}) This theorem is proved in essentially the same way in \gesperrt{Dirichlet-Dedekind}, Supplement II. Cf. also \gesperrt{Riemann's} mathematische Werke, 2. Auflage, Nr. XXX.
\endgroup
